% year: תשפ"ו
% type: בוחן
% semester: א
% official_solution: yes
% source: exams/בחנים/תשפ״ו, בוחן, 9141.pdf
% part: א
% number: 3
% points: 20
% categories: ivt

\begin{question}
הוכיחו שלמשוואה
\[
x^5+e^x=0
\]
יש פתרון ממשי, ומצאו קטע שאורכו לכל היותר 1 המכיל את הפתרון. נמקו.
\end{question}

\begin{hint}
הגדירו $f(x)=x^5+e^x$ ובדקו את הסימן שלה בנקודות שלמות קרובות ל-0.
\end{hint}

\begin{solution}
נגדיר $f(x)=x^5+e^x$. $f$ רציפה על $\R$ כסכום של פולינום ושל פונקציה מעריכית, שתיהן רציפות.

נחשב:
\[
f(0)=0^5+e^0=1>0,\qquad f(-1)=(-1)^5+e^{-1}=-1+\frac1e<0 ,
\]
כי $e>1$ ולכן $\frac1e<1$.

$f$ רציפה בקטע הסגור $[-1,0]$ ומחליפה בו סימן בקצוות. לפי משפט ערך הביניים (משפט בולצאנו) קיים $c\in(-1,0)$ כך ש-$f(c)=0$, כלומר $c^5+e^c=0$, ו-$c$ הוא פתרון ממשי של המשוואה.

\textbf{הקטע המבוקש:} $[-1,0]$, שאורכו 1.

(הערה: הפתרון יחיד — $f$ עולה ממש כסכום של שתי פונקציות עולות ממש $x^5$ ו-$e^x$ — ולכן "הפתרון" מוגדר היטב. מספרית $c\approx-0.845$.)
\end{solution}
